\documentclass[11pt]{article}

\usepackage{amsmath,amssymb,amstext,amsthm}
\usepackage{geometry}
\usepackage{fancyhdr}
\usepackage[pdftex]{graphicx}
\usepackage{xcolor}

\geometry{
  verbose,
  dvips,
  width=430pt, marginparsep=0pt, marginparwidth=0pt,
  top=90pt, headheight=12pt, headsep=20pt, footskip=30pt, bottom=80pt}

\setlength{\parskip}{\medskipamount}
\setlength{\parindent}{0pt}

\linespread{1.05}

\newtheorem{theorem}{Theorem}
\newtheorem{lemma}[theorem]{Lemma}
\newtheorem{proposition}[theorem]{Proposition}
\newtheorem{corollary}[theorem]{Corollary}
\newtheorem{conjecture}[theorem]{Conjecture}
\newtheorem{obs}{Observation}
\newtheorem{clm}{Claim}
\newtheorem{ques}{Question}
\newtheorem{problem}{Problem}

\newtheorem{ex}{Example}


\newcommand{\spc}{\hspace{0.08in}}
\newcommand{\vs}{\vspace*{.1in}}
\newcommand{\E}{\mathbb{E}}
\newcommand{\Prob}{\mathbb{P}}
\newcommand{\edit}[1]{\emph{\textcolor{blue}{#1}}}


\renewenvironment{proof}{{\noindent \textbf{\textit{Proof.}}}}
{\hfill $\Box$ \vspace*{0.1in}}
\newenvironment{proofc}{{\noindent \textit{Proof of Claim.}}}
{\hfill $\Box$}


\begin{document}

\begin{LARGE}\begin{center}\bf 13.3 Arc Length and Curvature\end{center}
\end{LARGE}

\vs\vs



\section{Warm Up}

\textbf{Question 1}: How can I find the length of a curve from calculus II?


\section{Motivation}




\section{Definitions \& Theorems}

\begin{enumerate}
\item  In two dimensions, using parametric equations $x = f(t)$ and $y=g(t)$, we get the length of an arc from $a \leq t \leq b$ to be

\begin{equation*}
    L = \int_a^b \sqrt{[f'(t)]^2+[g'(t)]^2}dt 
\end{equation*}

\item For a space curve, our length of a curve formula is really no different. We now have three dimensions. Suppose our space curve has the vector function $\langle f(t), g(t), h(t)\rangle$, the length of the space curve $L$ is as follows.

\begin{equation*}
  L=  \int_a^b \sqrt{[f'(t)]^2+[g'(t)]^2+[h'(t)]^2}dt = \int_a^b |r'(t)|dt.
\end{equation*}

\item We can now define an \textbf{arc length function} $s(t)$ which just takes our integral from an initial value $a$ to end end point $t$ which is variable.

\begin{equation*}
    s(t) = \int_a^t |r'(u)|du
\end{equation*}

Differentiating both sides we see that $ds/dt = |r'(t)|$ by the fundamental theorem of calculus. Since we now have $s$ as a function of $t$, i.e. a correspondence between $s$ and $t$, we can solve for $t$ and \textbf{re-parameterize} our vector function in terms of $s$.

\item If we want to consider how quickly a curve changes direction, it is natural to look at the rate at which the derivative changes with respect to the arc length at that given time. We call this \textbf{curvature} and it is defined as follows.

\begin{equation*}
    k = \left|\frac{dT}{ds}\right| = \left| \frac{dT/dt}{ds/dt} \right| = \left| \frac{T'(t)}{r'(t)}\right|.
\end{equation*}

\item Using our derivative rules (product rule) there is a very cool way to write curvature. The proof to the following theorem is omitted.

\begin{theorem}(10 in book)
    $k = \frac{|r'(t) \times r''(t)|}{|r'(t)|^3}$
\end{theorem}

\item Our unit tangent vector and arc length re-parameterization can help us out in another way. The entire purpose of our standard basis representation $\hat{i},\hat{j},\hat{k}$ is that every vector can be written as a linear combination of these. The only thing special about these though is that they are unit, and mutually orthogonal. They give us a global view of a space curve as if someone is looking at it from afar and can see the origin and everything as if its on a movie screen. What if instead there was a new basis system which was more beneficial from the perspective of the curve. If we think about a plane flying through a space curve, the pilot doesn't care much about i's, j's, and k's but rather their own view point. We therefore define the following vectors. We call them the \textbf{TNB-frames}

\begin{itemize}
    \item \textbf{T}: Unit Tangent Vector (Direction, nose of plane)
    \item \textbf{N}: Unit Normal Vector (Turning, tail-fin of plane)
    \item \textbf{B}: Unit Binormal Vector (Twist, wings of plane)
\end{itemize}

\begin{equation*}
  T(t) = \frac{r'(t)}{|r(t)|} \hspace{1cm} N(t) = \frac{T'(t)}{|T'(t)|} \hspace{1cm} B(t) = \frac{T(t) \times N(t)}{|T(t) \times N(t)|}=T(t)\times N(t).
\end{equation*}







   \end{enumerate}

\newpage


\section{Examples}

\begin{problem}
    Find the arc length of the vector function $\vec{r}(t) = \langle 2t, 3\sin(2t), 3\cos(2t) \rangle$ from $0\leq t \leq 2\pi$.
\end{problem}

\vspace{6cm}


\begin{problem}
    Find the arc length of the vector function $
    \vec{r}(t) = \frac{1}{3}t^3 \hat{i} + 4t \hat{j} + \sqrt{2}t^2 \hat{k}$ from $0\leq t\leq 2$.
\end{problem}

\vspace{6cm}

\begin{problem}
    Find the arc length function of the vector function from Problem 2.
\end{problem}

\newpage

\begin{problem}
    Reparameterize the following vector function in terms of arc length.

    \begin{equation*}
        \vec{r}(t) = \langle   e^t \cos(t), e^t \sin(t), e^t \rangle 
    \end{equation*}
\end{problem}


\vspace{8.5cm}


\begin{problem}
    Find the curvature of the vector function $\vec{r}(t) = \langle 4t, -t^2, 2t^3 \rangle$.
\end{problem}

\vspace{11cm}

\begin{problem}
    Find $T(t),N(t)$, and $B(t)$ if $r(t) = \cos(3t)\hat{i} + \sin(3t)\hat{j} + 5\hat{k}$.
\end{problem}

\newpage

\section{Scratch Work}

\end{document} 